% généré par scripts/04_generated.py
\begin{table}[tbp]\centering\footnotesize\setlength{\tabcolsep}{3.5pt}\caption{Contrôle des jours de la semaine (calendrier julien, recalcul ; l’année est celle de l’événement daté : naissance pour les \nnums{55} et 126 première mention).}\label{tab:jours}
\begin{tabular}{@{}r r l l l l l@{}}\toprule N° & Année & Mention & Quantième & Date & Calcul & Résultat\\\midrule
7 & 343 & \lat{die Martis} & \lat{III Nonas Iulias} & 5/7/343 & mardi & concordant \\
55 & 340 & \lat{die Veneris} & \lat{XI Kalendas Iulias} & 21/6/340 & samedi & écart d’un jour \\
118 & 367 & \lat{die Martis} & \lat{—} & — & — & quantième lacunaire \\
122 & 367 & \lat{die Martis} & \lat{Nonis Iuniis} & 5/6/367 & mardi & concordant \\
126 & 350 & \lat{die Saturnis} & \lat{III Idus Augustas} & 11/8/350 & samedi & concordant \\
126 & 368 & \lat{die Saturnis} & \lat{VII Idus Decembres} & 7/12/368 & dimanche & écart d’un jour \\
128 & 368 & \lat{die Martis} & \lat{XV Kalendas Decembres} & 17/11/368 & lundi & écart d’un jour \\
131 & 368 & \lat{die Veneris} & \lat{pridie Nonas Martias} & 6/3/368 & jeudi & écart d’un jour \\
169 & 373 & \lat{die Lunae} & \lat{XV Kalendas Apriles} & 18/3/373 & lundi & concordant \\
275 & 385 & \lat{die Lunae} & \lat{VI Idus Martia} & 10/3/385 & lundi & concordant \\
287 & 386 & \lat{die Solis} & \lat{X Kalendas Septembres} & 23/8/386 & dimanche & concordant \\
349 & 392 & \lat{die Veneris} & \lat{XV Kalendas Octobres} & 17/9/392 & vendredi & concordant \\
381 & 395 & \lat{die Saturni} & \lat{III Idus Maias} & 13/5/395 & dimanche & écart d’un jour \\
381 & 395 & \lat{die Lunae} & \lat{XII Kalendas Iunias} & 21/5/395 & lundi & concordant \\
416 & 397 & \lat{die Veneris} & \lat{tertium Kalendas Iulias} & 29/6/397 & lundi & discordant \\
428 & 398 & \lat{die Saturni} & \lat{XV Kalendas Ianuarias} & 18/12/398 & samedi & concordant \\
428 & 399 & \lat{die Saturni} & \lat{XV Kalendas Ianuarias} & 18/12/399 & dimanche & écart d’un jour \\
431 & 399 & \lat{die Veneris} & \lat{septimum Kalendas Septembres} & 26/8/399 & vendredi & concordant \\
462 & 402 & \lat{die Saturni} & \lat{—} & — & — & quantième lacunaire \\
\bottomrule\end{tabular}\end{table}
